\chapter{Categories}

Els exercicis següents completen els resolts al capítol corresponent de la part de Pràctica. Cada enunciat va acompanyat d'una indicació breu per si el lector s'encalla; convé, però, intentar-lo primer sense mirar-la. La marca \enquote{$\star$} assenyala un exercici de consolidació i \enquote{$\star\star$} un d'ampliació.

\section{Categories i axiomes}

\begin{exercici}
Comprova que un conjunt qualsevol $X$ dona lloc a una categoria \textbf{discreta}: els objectes són els elements de $X$ i els únics morfismes són les identitats. Verifica els axiomes. \hfill{\normalfont$\star$}
\begin{indicacio}
No hi ha cap morfisme entre objectes diferents; l'única composició possible en cada objecte és $\id_x\comp\id_x$.
\end{indicacio}
\end{exercici}

\begin{exercici}
Sigui $\cat{C}$ una categoria. Comprova que fixar un objecte $A$ i considerar només les fletxes $A\to A$, amb la composició de $\cat{C}$, dona un \textbf{monoide}. Quin és el seu element neutre? \hfill{\normalfont$\star$}
\begin{indicacio}
Els morfismes $A\to A$ es poden compondre entre ells i el resultat torna a anar d'$A$ a $A$. El neutre és $\id_A$.
\end{indicacio}
\end{exercici}

\begin{exercici}
Comprova que els espais topològics amb les aplicacions contínues formen una categoria $\cat{Top}$. Cal usar dues propietats de les aplicacions contínues: quines? \hfill{\normalfont$\star$}
\begin{indicacio}
La composició de contínues és contínua, i la identitat és contínua.
\end{indicacio}
\end{exercici}

\section{Fletxes, camins i diagrames commutatius}

\begin{exercici}
Escriu l'equació que expressa la commutativitat del triangle
\[
\begin{tikzpicture}[baseline=(m.center)]
  \matrix (m) [matrix of math nodes, column sep=3.3em, row sep=2.6em] {
    A & B \\
      & C \\
  };
  \draw[-{Stealth[scale=1.05]}] (m-1-1) -- node[above] {$f$} (m-1-2);
  \draw[-{Stealth[scale=1.05]}] (m-1-2) -- node[right] {$g$} (m-2-2);
  \draw[-{Stealth[scale=1.05]}] (m-1-1) -- node[below left] {$h$} (m-2-2);
\end{tikzpicture}
\]
i explica quins són els dos camins que es comparen. \hfill{\normalfont$\star$}
\begin{indicacio}
Tots dos camins comencen a $A$ i acaben a $C$; un passa per $B$ i l'altre és la fletxa directa.
\end{indicacio}
\end{exercici}

\begin{exercici}
Representa mitjançant un quadrat commutatiu l'equació
\[
q\comp u=v\comp p,
\]
on $u\colon A\to B$, $q\colon B\to D$, $p\colon A\to C$ i $v\colon C\to D$. \hfill{\normalfont$\star$}
\begin{indicacio}
Col·loca $A$ a l'angle superior esquerre i $D$ a l'inferior dret. Els dos membres de l'equació són dos camins d'$A$ a $D$.
\end{indicacio}
\end{exercici}

\begin{exercici}
Dona un exemple de quatre aplicacions entre conjunts que formin un quadrat que \emph{no} commuti. Indica explícitament un element sobre el qual les dues composicions donin resultats diferents. \hfill{\normalfont$\star$}
\begin{indicacio}
Pots treballar amb conjunts molt petits, per exemple $\{0,1\}$, i comparar els dos camins sobre l'element $0$.
\end{indicacio}
\end{exercici}

\section{Primers exemples}

\begin{exercici}
En un sistema deductiu $T$, suposem
\[
A\vdash_T B,
\qquad
B\vdash_T C,
\qquad
C\vdash_T A.
\]
Dibuixa el diagrama corresponent a $\cat{Prov}_T$ i comprova que $A$, $B$ i $C$ són isomorfs dos a dos. \hfill{\normalfont$\star$}
\begin{indicacio}
La transitivitat dona també $A\vdash_T C$, $B\vdash_T A$ i $C\vdash_T B$. En una categoria prima, dues fletxes en sentits oposats ja donen un isomorfisme.
\end{indicacio}
\end{exercici}

\begin{exercici}
Comprova directament que la composició de dos morfismes de grafs és un morfisme de grafs. Si $F\colon G\to H$ i $K\colon H\to L$, verifica les equacions d'origen i destí per a $K\comp F$. \hfill{\normalfont$\star$}
\begin{indicacio}
Per a l'origen, parteix de
\[
F_V\comp s_G=s_H\comp F_E,
\qquad
K_V\comp s_H=s_L\comp K_E,
\]
i compon les igualtats adequadament. Fes el mateix amb $t$.
\end{indicacio}
\end{exercici}

\begin{exercici}
Explica amb precisió quina informació té un graf dirigit que encara no té l'estructura d'una categoria. Després indica què afegeix la categoria lliure sobre aquest graf. \hfill{\normalfont$\star$}
\begin{indicacio}
Un graf té vèrtexs, arestes, origen i destí. Una categoria necessita, a més, identitats, composició i axiomes; la categoria lliure els genera mitjançant camins.
\end{indicacio}
\end{exercici}

\section{Isomorfismes}

\begin{exercici}
Comprova que a la categoria $\cat{Mon}$ dels monoides i els homomorfismes, un isomorfisme és exactament un homomorfisme bijectiu. Compara-ho amb el que pot passar a $\cat{Pos}$. \hfill{\normalfont$\star\star$}
\begin{indicacio}
Comprova que la inversa d'un homomorfisme bijectiu de monoides preserva l'operació i el neutre, de manera que torna a ser un homomorfisme. Per a $\cat{Pos}$, en canvi, la inversa d'una aplicació monòtona bijectiva pot no ser monòtona. Fes el contrast explícit amb un exemple: pren el conjunt $\{0,1\}$ amb l'ordre \emph{discret} (cap dels dos elements comparable amb l'altre) com a domini, i el mateix conjunt amb l'ordre $0<1$ com a codomini. La identitat subjacent és una aplicació monòtona (trivialment, perquè al domini no hi ha cap parella comparable) i bijectiva, però la seva inversa ---la identitat de l'ordre $0<1$ cap a l'ordre discret--- no és monòtona, ja que hauria d'enviar $0<1$ a dos elements incomparables. Per tant no és un isomorfisme a $\cat{Pos}$, tot i ser una bijecció monòtona.
\end{indicacio}
\end{exercici}

\begin{exercici}
Comprova que tot isomorfisme té secció. Després, prenent $A=\{0,1\}$, $B=\{\ast\}$ i $f$ l'única aplicació $A\to B$, dona una secció $s\colon B\to A$ explícita i comprova que $f$ no és isomorfisme. Observa que en aquest cas $f$ té \emph{dues} seccions distintes: quines? \hfill{\normalfont$\star$}
\begin{indicacio}
Tria $s(\ast)=0$ o $s(\ast)=1$. En tots dos casos $f\comp s=\id_B$, però $s\comp f\neq\id_A$. Les dues tries donen les dues seccions.
\end{indicacio}
\end{exercici}

\begin{exercici}
A la categoria $\cat{Set}$, exhibeix morfismes que il·lustrin els tres casos possibles pel que fa a l'existència de seccions: un morfisme sense cap secció, un amb una única secció i un amb més d'una. \hfill{\normalfont$\star\star$}
\begin{indicacio}
Per a \emph{cap} secció, cerca un morfisme no exhaustiu, com $\emptyset\to\{\ast\}$ o una injecció no exhaustiva. Per a \emph{una única}, un morfisme exhaustiu on cada element del codomini tingui una sola antiimatge (una bijecció). Per a \emph{diverses}, un morfisme exhaustiu amb alguna fibra de més d'un element: cada elecció d'antiimatges dona una secció distinta. Recorda que a $\cat{Set}$ triar una secció d'una aplicació exhaustiva és triar un element a cada fibra.
\end{indicacio}
\end{exercici}

\begin{exercici}
Comprova que si $g\comp f$ és un isomorfisme, no cal que ho siguin ni $f$ ni $g$. Busca'n un exemple senzill a $\cat{Set}$ i representa'l amb fletxes. \hfill{\normalfont$\star$}
\begin{indicacio}
Pensa en una injecció seguida d'una projecció entre conjunts de mides diferents.
\end{indicacio}
\end{exercici}

\begin{exercici}
Demostra la caracterització dual de l'exercici resolt a la part de pràctica: un morfisme $f$ és un isomorfisme si i només si és un epimorfisme escindit (té secció) i un monomorfisme. \hfill{\normalfont$\star$}
\begin{indicacio}
És el dual de \enquote{mono escindit i epi $\Rightarrow$ iso}: parteix d'una secció $s$ amb $f\comp s=\id_B$ i fes servir que $f$ és mono per veure que $s\comp f=\id_A$.
\end{indicacio}
\end{exercici}

\begin{exercici}
Considera la categoria $\cat{2}=(0\to 1)$ i la seva única fletxa no trivial $u\colon 0\to 1$. Comprova que $u$ és alhora un monomorfisme i un epimorfisme (un bimorfisme), però que \emph{no} és un isomorfisme. Quin fet elemental de $\cat{2}$ fa que les cancel·lacions siguin trivials? \hfill{\normalfont$\star$}
\begin{indicacio}
A $\cat{2}$ els únics morfismes són $\id_0$, $\id_1$ i $u$. Per veure que $u$ és mono i epi, comprova que amb tan pocs morfismes les úniques parelles $g,h$ que compleixen les hipòtesis de cancel·lació ja són iguals de partida. Per veure que no és iso, observa que no hi ha cap fletxa $1\to 0$, de manera que $u$ no pot tenir invers.
\end{indicacio}
\end{exercici}

\begin{exercici}
Estudia quin factor hereta la propietat en una composició. Demostra que si $g\comp f$ és un monomorfisme, aleshores $f$ és un monomorfisme; i que, dualment, si $g\comp f$ és un epimorfisme, aleshores $g$ és un epimorfisme. Comprova amb un contraexemple que, en el primer cas, $g$ no cal que sigui mono, i en el segon, $f$ no cal que sigui epi. \hfill{\normalfont$\star\star$}
\begin{indicacio}
Per al mono: de $f\comp u=f\comp v$, compon per l'esquerra amb $g$ per obtenir $(g\comp f)\comp u=(g\comp f)\comp v$ i aplica que $g\comp f$ és mono. El cas de l'epi és el dual: compon per la dreta. Observa l'asimetria: d'un mono $g\comp f$ hereta el factor de la \emph{dreta} ($f$), i d'un epi, el de l'\emph{esquerra} ($g$). Per als contraexemples, pensa en una injecció i una projecció a $\cat{Set}$.
\end{indicacio}
\end{exercici}

\section{Construccions sobre categories}

\begin{exercici}
Descriu la categoria producte $\cat{3}\times\cat{3}$, on $\cat{3}$ és la categoria associada a l'ordre $0<1<2$. Quants objectes i quants morfismes té? Dibuixa les fletxes bàsiques i explica quines altres fletxes apareixen per composició. \hfill{\normalfont$\star\star$}
\begin{indicacio}
$\cat{3}$ té tres objectes i sis morfismes. En el producte, els objectes i els morfismes són parelles.
\end{indicacio}
\end{exercici}

\begin{exercici}
Compara directament $(\cat{C}\times\cat{D})\op$ i $\cat{C}\op\times\cat{D}\op$. Comprova que tenen els mateixos objectes i que els seus morfismes, identitats i composicions s'identifiquen component a component. \hfill{\normalfont$\star$}
\begin{indicacio}
Un morfisme $(A,X)\to(B,Y)$ a $(\cat{C}\times\cat{D})\op$ és, per definició, un morfisme $(B,Y)\to(A,X)$ a $\cat{C}\times\cat{D}$. Escriu què significa això en cadascuna de les dues components.
\end{indicacio}
\end{exercici}

\begin{exercici}
Per a cadascuna de les subcategories següents de $\cat{Set}$, digues si és \textbf{plena}, \textbf{ampla}, totes dues coses o cap: (a)~$\cat{FinSet}$, amb els conjunts finits per objectes i totes les aplicacions entre ells; (b)~la subcategoria amb tots els conjunts per objectes i només les aplicacions injectives; (c)~la subcategoria amb tots els conjunts per objectes i totes les aplicacions (és a dir, $\cat{Set}$ mateixa); (d)~la subcategoria amb els conjunts finits per objectes i només les aplicacions injectives entre ells. \hfill{\normalfont$\star$}
\begin{indicacio}
Recorda que \emph{plena} vol dir que es conserven \emph{tots} els morfismes entre els objectes escollits (la subcategoria queda determinada només pels objectes), i \emph{ampla} vol dir que es conserven \emph{tots} els objectes. Comprova cada cas contra aquestes dues condicions per separat: poden fallar independentment. En particular, mira si (d) és plena o ampla, i observa que una subcategoria pot no ser ni l'una ni l'altra.
\end{indicacio}
\end{exercici}

\section{Grafs i categories lliures}

\begin{exercici}
Classifica, \textbf{fins a isomorfisme}, els grafs dirigits que generen una categoria lliure amb exactament \textbf{quatre} morfismes. Recorda que les identitats compten com a morfismes. \hfill{\normalfont$\star\star$}
\begin{indicacio}
Els morfismes de la categoria lliure són els camins dirigits del graf, inclosos els camins buits ---un per vèrtex--- que fan d'identitats. Un cicle dirigit produiria infinits camins, de manera que el graf ha de ser acíclic; aleshores el nombre total de morfismes és finit. La classificació té exactament tres casos fins a isomorfisme:
\begin{enumerate}
  \item quatre vèrtexs i cap aresta;
  \item tres vèrtexs i una única aresta entre dos vèrtexs diferents, amb el tercer aïllat;
  \item dos vèrtexs i dues arestes paral·leles en el mateix sentit.
\end{enumerate}
En tots tres no hi ha cap camí de longitud superior a $1$, i el total de morfismes és $|V_G|+|E_G|=4$. Per veure que no n'hi ha d'altres: sense cicles, quatre morfismes deixen com a màxim quatre vèrtexs (cada vèrtex ja aporta la seva identitat). Amb un sol vèrtex només hi ha la identitat, i qualsevol aresta hi seria un bucle, que dona infinits morfismes. Amb dos vèrtexs, les dues identitats deixen espai per a dues arestes més, que han de ser paral·leles en el mateix sentit (una en cada sentit formaria un cicle). Amb tres vèrtexs, les tres identitats deixen espai per a una sola aresta. Amb quatre, les quatre identitats ja exhaureixen el compte i no hi cap cap aresta.
\end{indicacio}
\end{exercici}

\begin{exercici}
Considera el graf
\[
A\xrightarrow{f}B\xrightarrow{g}C.
\]
Llista tots els morfismes de la categoria lliure que genera, dibuixa també la fletxa composta i comprova que hi ha exactament sis morfismes. \hfill{\normalfont$\star$}
\begin{indicacio}
Compta els tres camins buits, les dues arestes $f$ i $g$, i el camí compost de longitud dos.
\end{indicacio}
\end{exercici}

\begin{exercici}
Per què un graf amb un cicle, per exemple una aresta $a\colon A\to A$, genera sempre una categoria lliure amb infinits morfismes? \hfill{\normalfont$\star$}
\begin{indicacio}
El cicle es pot recórrer tantes vegades com es vulgui: $a,a^2,a^3,\ldots$, i cada repetició dona un camí diferent.
\end{indicacio}
\end{exercici}

\section{Qüestions de mida}

\begin{exercici}
Comprova que un preordre, vist com a categoria, és una categoria \textbf{petita} sempre que el conjunt subjacent $P$ sigui un conjunt. Quants morfismes hi ha, com a màxim, entre dos objectes? \hfill{\normalfont$\star$}
\begin{indicacio}
Entre dos objectes hi ha zero o un morfisme; la col·lecció total de morfismes es pot identificar amb un subconjunt de $P\times P$.
\end{indicacio}
\end{exercici}

\begin{exercici}
Explica per què $\cat{Set}$ és localment petita. És a dir, per què, donats dos conjunts $A$ i $B$, la col·lecció d'aplicacions $A\to B$ és un conjunt. \hfill{\normalfont$\star$}
\begin{indicacio}
Les aplicacions $A\to B$ es poden identificar amb certs subconjunts de $A\times B$; per tant formen un subconjunt de $\mathcal P(A\times B)$.
\end{indicacio}
\end{exercici}

\begin{exercici}
Explica per què $\cat{Graph}$ és localment petita: donats dos grafs dirigits $G$ i $H$, per què els morfismes de grafs $G\to H$ formen un conjunt? \hfill{\normalfont$\star$}
\begin{indicacio}
Un morfisme de grafs és una parella d'aplicacions $V_G\to V_H$ i $E_G\to E_H$ que satisfà dues condicions. Aquestes parelles formen un subconjunt d'un producte de dos conjunts d'aplicacions.
\end{indicacio}
\end{exercici}
